% ============================================================
%  CHAPTER 5: WORK, ENERGY AND POWER
%  The Physics Codex: A Complete University Foundation
%  Korede Samuel Omotosho (Falcon)
%  Aurikrex Learning Systems, 2026
%
%  File: chapters/part1/ch05_work_energy_power.tex
% ============================================================

\chapter{Work, Energy and Power}
\label{ch:work_energy_power}

\chapterquote{Energy cannot be created or destroyed;
it can only be changed from one form to another.}{Albert Einstein}

\begin{objectivebox}
By the end of this chapter, you should be able to:
\begin{enumerate}[noitemsep]
  \item Define work and calculate it for constant and
  variable forces
  \item State and apply the work-energy theorem
  \item Define kinetic and potential energy and apply
  conservation of mechanical energy
  \item Calculate power and relate it to force and
  velocity
  \item Define efficiency and apply it to machines and
  energy transfers
  \item Identify different forms of energy and their
  interconversions
  \item Solve problems involving springs and elastic
  potential energy
\end{enumerate}
\end{objectivebox}

\vspace{0.4cm}

\minitoc

\vspace{0.5cm}

% ============================================================
\section{The Concept of Work}
% ============================================================

In everyday language, ``work'' means any effortful
activity --- studying, walking, thinking. In Physics,
work has a precise and narrower meaning. You can push
against a wall with all your strength for an hour and
do zero work in the Physics sense, because the wall
did not move.

\begin{definitionbox}[Work Done by a Constant Force]
The \textbf{work} $W$ done by a constant force $\vect{F}$
when an object moves through displacement $\vect{s}$ is:
\[
  W = Fs\cos\theta
\]
where $\theta$ is the angle between the force and the
displacement. Work is a \textbf{scalar}. SI unit:
joule (J) where $1\ \text{J} = 1\ \text{N\,m}$.
\end{definitionbox}

\begin{diagbox}[Work Done at an Angle]
\begin{center}
\begin{tikzpicture}[scale=1.0]
  % Ground
  \fill[gray!20] (0,-0.3) rectangle (8,0);
  \draw[thick] (0,0) -- (8,0);

  % Object
  \fill[codexblue!30] (1,0) rectangle (2.5,0.8);
  \fill[codexblue!30] (5.5,0) rectangle (7,0.8);
  \node[font=\small] at (1.75,0.4) {$m$};
  \draw[->, dashed, gray, thick] (2.5,0.4) -- (5.5,0.4);
  \node[above, gray] at (4,0.4) {$s$};

  % Force arrow
  \draw[->, very thick, codexred]
    (1.75,0.8) -- ++(3,1.5)
    node[right]{$\vect{F}$};

  % Angle
  \draw[->, dashed, gray] (1.75,0.8) -- ++(3,0);
  \draw[codexgold] (2.55,0.8) arc(0:26.6:0.8);
  \node[codexgold, anchor=east] at (2.65,1.0) {$\theta$};

  % Horizontal component
  \draw[->, thick, codexblue, dashed]
    (1.75,0.8) -- ++(3,0)
    node[pos=0.7, above]{$F\cos\theta$};

  % Work formula
  \node[align=center, font=\small] at (4,-0.8)
    {Work $= F\cos\theta \times s$\\
    (only the component along displacement does work)};
\end{tikzpicture}
\end{center}
\end{diagbox}

\begin{keybox}[When is Work Zero?]
Work done is zero when:
\begin{itemize}[noitemsep]
  \item The force is zero ($F = 0$)
  \item The displacement is zero ($s = 0$) --- no
  movement
  \item The force is perpendicular to displacement
  ($\theta = 90°$, so $\cos 90° = 0$)
\end{itemize}
Example: A waiter carrying a tray horizontally does
\textbf{zero work} on the tray with the upward
force of his hand (force is vertical, displacement
is horizontal, $\theta = 90°$). The effort he feels
is from maintaining the force, not from doing work
in the Physics sense.
\end{keybox}

\begin{examplebox}[5.1]
\stepcounter{workedexample}
A person drags a $20\ \text{kg}$ box across a floor
with a rope at $35°$ above the horizontal. The tension
in the rope is $80\ \text{N}$. The box moves $5.0\ \text{m}$.
Find:\\
(a) the work done by the tension;\\
(b) the work done by gravity;\\
(c) the work done by the normal reaction.
\end{examplebox}

\begin{solutionbox}
\textbf{(a)} Work by tension:
\[
  W_T = F\cos\theta \cdot s
  = 80 \times \cos 35° \times 5.0
  = 80 \times 0.819 \times 5.0
  = 327.6 \approx 328\ \text{J}
\]

\textbf{(b)} Work by gravity: Gravity acts downward
($-y$), displacement is horizontal. The angle between
gravity and displacement is $90°$.
\[
  W_g = mg\cos 90° \times s = 0\ \text{J}
\]

\textbf{(c)} Work by normal reaction: Normal reaction
acts upward ($+y$), displacement is horizontal.
Angle is $90°$.
\[
  W_N = N\cos 90° \times s = 0\ \text{J}
\]

Only the applied force does work in this problem.
\end{solutionbox}

% ============================================================
\section{Kinetic Energy}
% ============================================================

\begin{definitionbox}[Kinetic Energy]
The \textbf{kinetic energy} (KE) of a body is the energy
it possesses due to its motion:
\[
  KE = \frac{1}{2}mv^2
\]
KE is a \textbf{scalar}, always positive or zero.
SI unit: joule (J).
\end{definitionbox}

\begin{lawbox}[Work-Energy Theorem]
The \textbf{net work} done on a body equals the change
in its kinetic energy:
\[
  W_{net} = \Delta KE = \frac{1}{2}mv^2 - \frac{1}{2}mu^2
\]
This is one of the most powerful results in classical
mechanics. It connects forces and motion directly,
bypassing the need to know acceleration explicitly.
\end{lawbox}

\textbf{Derivation:}
From Newton's second law with constant force:
$F = ma$, and from kinematics: $v^2 = u^2 + 2as$,
so $as = (v^2-u^2)/2$. Then:
\[
  W = Fs = (ma)s = m(as) = m\cdot\frac{v^2-u^2}{2}
  = \tfrac{1}{2}mv^2 - \tfrac{1}{2}mu^2 = \Delta KE
\]

\begin{examplebox}[5.2]
\stepcounter{workedexample}
A $1000\ \text{kg}$ car accelerates from
$10\ \text{m\,s}^{-1}$ to $25\ \text{m\,s}^{-1}$
while a net force acts on it over $150\ \text{m}$.\\
(a) Find the net work done.\\
(b) Find the net force.
\end{examplebox}

\begin{solutionbox}
\textbf{(a)} Net work = change in KE:
\[
  W_{net} = \tfrac{1}{2}(1000)(25^2)
  - \tfrac{1}{2}(1000)(10^2)
\]
\[
  = 312\,500 - 50\,000 = 262\,500\ \text{J}
  \approx 263\ \text{kJ}
\]

\textbf{(b)} From $W = Fs$:
\[
  F = \frac{W_{net}}{s}
  = \frac{262\,500}{150} = 1\,750\ \text{N}
\]
\end{solutionbox}

% ============================================================
\section{Gravitational Potential Energy}
% ============================================================

\begin{definitionbox}[Gravitational Potential Energy]
The \textbf{gravitational potential energy} (GPE) of
a body at height $h$ above a reference level is:
\[
  GPE = mgh
\]
It is the work that must be done against gravity to
raise the body from the reference level to height $h$.
GPE depends on the choice of reference level (usually
the ground or lowest point in the problem).
\end{definitionbox}

\begin{lawbox}[Conservation of Mechanical Energy]
In the absence of friction and air resistance, the
total mechanical energy (KE + GPE) of a system is
conserved:
\[
  E_{mechanical} = KE + GPE
  = \frac{1}{2}mv^2 + mgh = \text{constant}
\]
\[
  \frac{1}{2}mu^2 + mgh_1
  = \frac{1}{2}mv^2 + mgh_2
\]
What is lost in KE is gained in GPE, and vice versa.
\end{lawbox}

\begin{examplebox}[5.3]
\stepcounter{workedexample}
A ball is released from rest at the top of a smooth
ramp $3.2\ \text{m}$ above the ground.
Find the speed of the ball at the bottom of the ramp.
($g = 10\ \text{m\,s}^{-2}$)
\end{examplebox}

\begin{solutionbox}
At the top: $KE = 0$, $GPE = mgh = mg(3.2)$\\
At the bottom: $KE = \frac{1}{2}mv^2$, $GPE = 0$

By conservation of energy:
\[
  mgh = \frac{1}{2}mv^2
\]
\[
  v^2 = 2gh = 2 \times 10 \times 3.2 = 64
\]
\[
  v = 8.0\ \text{m\,s}^{-1}
\]

Note: the mass cancels. The speed at the bottom is
independent of the mass --- exactly as Galileo showed
for free fall.
\end{solutionbox}

\begin{examplebox}[5.4]
\stepcounter{workedexample}
A $500\ \text{g}$ ball is thrown vertically upward
with speed $14\ \text{m\,s}^{-1}$. Using energy
conservation, find:\\
(a) the maximum height;\\
(b) the speed when it is at height $5.0\ \text{m}$.
($g = 10\ \text{m\,s}^{-2}$)
\end{examplebox}

\begin{solutionbox}
\textbf{(a)} At maximum height, $v = 0$:
\[
  \frac{1}{2}mv^2 + mgh_{max}
  = \frac{1}{2}mu^2 + 0
\]
\[
  0 + mgh_{max} = \frac{1}{2}mu^2
\]
\[
  h_{max} = \frac{u^2}{2g}
  = \frac{14^2}{2 \times 10} = \frac{196}{20}
  = 9.8\ \text{m}
\]

\textbf{(b)} At $h = 5.0\ \text{m}$:
\[
  \frac{1}{2}mv^2 = \frac{1}{2}mu^2 - mgh
\]
\[
  v^2 = u^2 - 2gh = 196 - 2(10)(5.0) = 96
\]
\[
  v = \sqrt{96} = 9.8\ \text{m\,s}^{-1}
\]
\end{solutionbox}

% ============================================================
\section{Elastic Potential Energy}
% ============================================================

\begin{definitionbox}[Elastic Potential Energy]
When a spring (or any elastic material) is compressed
or stretched by extension $x$ from its natural length,
energy is stored as \textbf{elastic potential energy}:
\[
  E_{elastic} = \frac{1}{2}kx^2
\]
where $k$ is the spring constant (N\,m$^{-1}$).
This stored energy equals the work done in compressing
or stretching the spring.
\end{definitionbox}

\begin{diagbox}[Work Done in Stretching a Spring]
\begin{center}
\begin{tikzpicture}
\begin{axis}[
  width=8cm, height=5cm,
  xlabel={Extension $x$ (m)},
  ylabel={Force $F$ (N)},
  xmin=0, xmax=5, ymin=0, ymax=30,
  grid=major, grid style={gray!20},
]
  % F = kx line (k=5)
  \addplot[codexblue, very thick, domain=0:5]{5*x};
  \addlegendentry{$F = kx$}

  % Shaded area = work done = elastic PE
  \addplot[fill=codexblue!20, opacity=0.6, domain=0:5]
    {5*x} \closedcycle;

  % Annotation
  \node[codexblue, font=\small, fill=white, inner sep=2pt]
    at (axis cs:1.7,22.5)
    {Area $= \frac{1}{2}kx^2 = E_{elastic}$};
\end{axis}
\end{tikzpicture}
\end{center}
\small The area under a force-extension graph equals
the work done in stretching the spring (and the elastic PE stored).
\end{diagbox}

% ============================================================
\section{Power}
% ============================================================

\begin{definitionbox}[Power]
\textbf{Power} is the rate of doing work, or the rate
of energy transfer:
\[
  P_{\mathrm{avg}} = \frac{\Delta W}{\Delta t}
  = \frac{\Delta E}{\Delta t}
\]
Instantaneous power is
\[
  P = \frac{dW}{dt} = \frac{dE}{dt}
  = \vect{F}\cdot\vect{v} = Fv\cos\theta
\]
If the force is parallel to the velocity:
$P = Fv$.

SI unit: watt (W) where $1\ \text{W} = 1\ \text{J\,s}^{-1}$.
Also used: kilowatt (kW = $10^3$ W),
megawatt (MW = $10^6$ W), horsepower (1 hp $\approx$ 746 W).
\end{definitionbox}

\begin{historybox}[The Watt and James Watt]
The unit of power is named after James Watt (1736--1819),
the Scottish engineer who dramatically improved the steam
engine and powered the Industrial Revolution. Watt
himself introduced the unit ``horsepower'' to describe
his engines' output in terms that horse-owners could
relate to. He estimated that a horse could do
33,000 foot-pounds of work per minute, and used this
to market his engines. One horsepower is approximately
746 watts.
\end{historybox}

\begin{examplebox}[5.5]
\stepcounter{workedexample}
A motor lifts a $200\ \text{kg}$ load through
$15\ \text{m}$ in $10\ \text{s}$.\\
(a) Find the minimum power of the motor.\\
(b) If the motor is $80\%$ efficient, what is the
actual power input?
($g = 10\ \text{m\,s}^{-2}$)
\end{examplebox}

\begin{solutionbox}
\textbf{(a)} Work done against gravity:
\[
  W = mgh = 200 \times 10 \times 15 = 30\,000\ \text{J}
\]
Minimum power (assuming 100\% efficiency):
\[
  P_{out} = \frac{W}{t} = \frac{30\,000}{10}
  = 3\,000\ \text{W} = 3.0\ \text{kW}
\]

\textbf{(b)} Efficiency $\eta = P_{out}/P_{in}$:
\[
  P_{in} = \frac{P_{out}}{\eta}
  = \frac{3000}{0.80} = 3\,750\ \text{W} = 3.75\ \text{kW}
\]
\end{solutionbox}

\begin{examplebox}[5.6]
\stepcounter{workedexample}
A car engine produces a driving force of $4000\ \text{N}$
when the car moves at $30\ \text{m\,s}^{-1}$ at constant
speed. Calculate the power output of the engine.
\end{examplebox}

\begin{solutionbox}
At constant speed, the driving force equals the
resistance. Power:
\[
  P = Fv = 4000 \times 30 = 120\,000\ \text{W}
  = 120\ \text{kW}
\]
\end{solutionbox}

% ============================================================
\section{Efficiency}
% ============================================================

\begin{definitionbox}[Efficiency]
The \textbf{efficiency} $\eta$ of an energy transfer
or machine is the ratio of useful output to total input,
expressed as a percentage:
\[
  \eta = \frac{\text{useful output energy}}
  {\text{total input energy}} \times 100\%
  = \frac{P_{out}}{P_{in}} \times 100\%
\]
Efficiency is always less than 100\% in practice,
because some energy is always dissipated (usually
as heat and sound) due to friction and other losses.
\end{definitionbox}

\begin{realworldbox}[Generator Efficiency in Nigeria]
A typical petrol generator commonly used in Nigerian
homes has an efficiency of around 20--30\%. This means
that for every $1\,000\ \text{J}$ of chemical energy
burned in petrol, only $200$--$300\ \text{J}$ is
converted to electrical energy. The remaining 70--80\%
is lost as heat in the engine and exhaust. This is why
generators feel very hot when running --- they are
radiating enormous amounts of wasted energy.
\end{realworldbox}

% ============================================================
\section{Forms of Energy and Energy Conservation}
% ============================================================

Energy exists in many forms. The principle of
conservation of energy states that energy can never
be created or destroyed --- only converted from one
form to another.

\begin{defbox}[Forms of Energy]
\begin{itemize}[noitemsep]
  \item \textbf{Kinetic energy:} energy of motion
  \item \textbf{Gravitational potential energy:}
  energy due to position in gravitational field
  \item \textbf{Elastic potential energy:} energy
  stored in deformed elastic materials
  \item \textbf{Chemical energy:} stored in chemical
  bonds (food, fuel, batteries)
  \item \textbf{Nuclear energy:} stored in atomic
  nuclei (fission, fusion)
  \item \textbf{Thermal (heat) energy:} internal
  energy due to random molecular motion
  \item \textbf{Electrical energy:} energy of moving
  charges
  \item \textbf{Radiant (light/EM) energy:}
  energy carried by electromagnetic waves
  \item \textbf{Sound energy:} mechanical wave energy
  in a medium
\end{itemize}
\end{defbox}

\begin{realworldbox}[Energy Chains in a Hydroelectric Dam]
At Kainji Dam on the Niger River (Nigeria's first
hydroelectric power station, commissioned 1968),
the energy conversion chain is:
\[
  \underbrace{\text{Gravitational PE}}_{\text{water at height}}
  \to \underbrace{\text{KE}}_{\text{falling water}}
  \to \underbrace{\text{Kinetic (turbine)}}_{\text{spinning}}
  \to \underbrace{\text{Electrical}}_{\text{generator}}
  \to \underbrace{\text{Light, heat, work}}_{\text{in homes}}
\]
Each conversion loses some energy as heat, which is why
the efficiency of the whole chain is well below 100\%.
\end{realworldbox}

% ============================================================
\section{Energy Dissipation by Friction}
% ============================================================

When friction acts, kinetic energy is converted to heat.
The work done against friction is:
\[
  W_{friction} = f \times d = \mu mg\cos\theta \times d
\]

By the work-energy theorem, including friction:
\[
  KE_{final} = KE_{initial} + W_{net forces}
  - W_{friction}
\]
Or equivalently:
\[
  \frac{1}{2}mv_2^2 = \frac{1}{2}mv_1^2
  + mgh - W_{friction}
\]

\begin{examplebox}[5.7]
\stepcounter{workedexample}
A $2.0\ \text{kg}$ block slides $4.0\ \text{m}$ down
a rough incline of $30°$ from rest. The coefficient
of kinetic friction is $0.25$. Find the speed at the
bottom. ($g = 10\ \text{m\,s}^{-2}$)
\end{examplebox}

\begin{solutionbox}
Height descended: $h = 4.0\sin 30° = 2.0\ \text{m}$\\
Normal force: $N = mg\cos 30° = 2.0(10)(0.866) = 17.3\ \text{N}$\\
Friction force: $f = \mu N = 0.25 \times 17.3 = 4.33\ \text{N}$\\
Work done by friction: $W_f = -4.33 \times 4.0 = -17.3\ \text{J}$

Energy conservation with friction:
\[
  \frac{1}{2}mv^2 = mgh - W_f
\]
\[
  \frac{1}{2}(2.0)v^2 = (2.0)(10)(2.0) - 17.3
  = 40 - 17.3 = 22.7\ \text{J}
\]
\[
  v^2 = \frac{22.7}{1.0} = 22.7
  \implies v = 4.8\ \text{m\,s}^{-1}
\]
\end{solutionbox}

% ============================================================
\section{Chapter Summary}
% ============================================================

\begin{summarybox}
\textbf{Work:} $W = Fs\cos\theta$ (scalar, J).
Zero when force $\perp$ displacement.

\textbf{Work-energy theorem:}
$W_{net} = \Delta KE = \frac{1}{2}mv^2 - \frac{1}{2}mu^2$

\textbf{KE:} $\frac{1}{2}mv^2$ \quad
\textbf{GPE:} $mgh$ \quad
\textbf{Elastic PE:} $\frac{1}{2}kx^2$

\textbf{Conservation of mechanical energy}
(no friction):
$\frac{1}{2}mv^2 + mgh = \text{constant}$

\textbf{Power:}
$P = W/t = Fv$ (watts, W)

\textbf{Efficiency:}
$\eta = (P_{out}/P_{in}) \times 100\%$

\textbf{Energy} cannot be created or destroyed,
only converted. Friction converts KE to heat.
\end{summarybox}

% ============================================================
\newpage
\begin{exercisebox}[5]

\subsection*{Section A: Multiple Choice}

\textbf{A1.} A force of $50\ \text{N}$ acts at $60°$
to the direction of motion of an object. The work done
when the object moves $8.0\ \text{m}$ is:

\mcoption{A}{$200\ \text{J}$}
\mcoption{B}{$346\ \text{J}$}
\mcoption{C}{$400\ \text{J}$}
\mcoption{D}{$692\ \text{J}$}
\ansref{Ch5-A1}

\vspace{0.4cm}

\textbf{A2.} A body falls freely from rest and reaches
speed $v$ after time $t$. Its kinetic energy at this
point is:

\mcoption{A}{$mgv$}
\mcoption{B}{$\frac{1}{2}mg^2t^2$}
\mcoption{C}{$\frac{1}{2}mv^2$}
\mcoption{D}{$mgt$}
\ansref{Ch5-A2}

\vspace{0.4cm}

\textbf{A3.} A spring with spring constant
$400\ \text{N\,m}^{-1}$ is compressed by $5.0\ \text{cm}$.
The elastic potential energy stored is:

\mcoption{A}{$0.25\ \text{J}$}
\mcoption{B}{$0.50\ \text{J}$}
\mcoption{C}{$10\ \text{J}$}
\mcoption{D}{$20\ \text{J}$}
\ansref{Ch5-A3}

\vspace{0.4cm}

\textbf{A4.} A pump lifts water at a rate of
$50\ \text{kg\,s}^{-1}$ through a height of $10\ \text{m}$.
The minimum power of the pump is
($g = 10\ \text{m\,s}^{-2}$):

\mcoption{A}{$500\ \text{W}$}
\mcoption{B}{$5\,000\ \text{W}$}
\mcoption{C}{$50\,000\ \text{W}$}
\mcoption{D}{$500\,000\ \text{W}$}
\ansref{Ch5-A4}

\vspace{0.4cm}

\textbf{A5.} A car engine exerts a force of
$3\,000\ \text{N}$ against a resistance of
$1\,000\ \text{N}$. At speed $20\ \text{m\,s}^{-1}$,
the power of the engine is:

\mcoption{A}{$20\,000\ \text{W}$}
\mcoption{B}{$40\,000\ \text{W}$}
\mcoption{C}{$60\,000\ \text{W}$}
\mcoption{D}{$80\,000\ \text{W}$}
\ansref{Ch5-A5}

\vspace{0.4cm}

\textbf{A6.} A machine has input power $2\ \text{kW}$
and output power $1.5\ \text{kW}$. Its efficiency is:

\mcoption{A}{$25\%$}
\mcoption{B}{$50\%$}
\mcoption{C}{$75\%$}
\mcoption{D}{$133\%$}
\ansref{Ch5-A6}

\vspace{0.4cm}

\textbf{A7.} Which of the following does zero work
on a planet in circular orbit around the Sun?

\mcoption{A}{Gravity}
\mcoption{B}{The planet's velocity}
\mcoption{C}{Atmospheric drag}
\mcoption{D}{Gravity does zero work because it acts
perpendicular to the orbital displacement at all times}
\ansref{Ch5-A7}

\vspace{0.4cm}

\textbf{A8.} A ball of mass $m$ rolls from the top of
a smooth hill of height $H$ with initial speed $v_0$.
Its speed at the bottom is:

\mcoption{A}{$\sqrt{v_0^2 + gH}$}
\mcoption{B}{$\sqrt{v_0^2 + 2gH}$}
\mcoption{C}{$\sqrt{2gH}$}
\mcoption{D}{$v_0 + \sqrt{2gH}$}
\ansref{Ch5-A8}

\vspace{0.4cm}

\textbf{A9.} The kinetic energy of an object is
doubled. Its speed is multiplied by a factor of:

\mcoption{A}{$\sqrt{2}$}
\mcoption{B}{$2$}
\mcoption{C}{$4$}
\mcoption{D}{$\sqrt{4} = 2$, same as B}
\ansref{Ch5-A9}

\vspace{0.4cm}

\textbf{A10.} A $70\ \text{kg}$ student climbs a
staircase of vertical height $12\ \text{m}$ in
$18\ \text{s}$. The power output of the student is
($g = 10\ \text{m\,s}^{-2}$):

\mcoption{A}{$350\ \text{W}$}
\mcoption{B}{$467\ \text{W}$}
\mcoption{C}{$580\ \text{W}$}
\mcoption{D}{$840\ \text{W}$}
\ansref{Ch5-A10}

\vspace{0.6cm}

\subsection*{Section B: Theory and Calculations}

\textbf{B1.} A $3.0\ \text{kg}$ block is pulled
$6.0\ \text{m}$ along a rough horizontal surface
by a horizontal force of $25\ \text{N}$. The
coefficient of kinetic friction is $0.30$.
($g = 10\ \text{m\,s}^{-2}$)\\[4pt]
(a) Find the work done by the applied force.\\
(b) Find the work done by friction.\\
(c) Find the net work done on the block.\\
(d) If the block started from rest, find its speed
after $6.0\ \text{m}$ using the work-energy theorem.\\
(e) Verify your answer using the equations of motion.
\hfill\ansref{Ch5-B1}

\vspace{0.4cm}

\textbf{B2.} A ball of mass $0.50\ \text{kg}$ is
attached to a string and swings in a vertical circle
of radius $1.2\ \text{m}$. At the lowest point of
the circle it has speed $6.0\ \text{m\,s}^{-1}$.
($g = 10\ \text{m\,s}^{-2}$)\\[4pt]
(a) Calculate the KE and GPE at the lowest point
(take lowest point as reference level).\\
(b) Using energy conservation, find the speed at
the highest point.\\
(c) What is the minimum speed at the highest point
for the string to remain taut? (Hint: at minimum
speed, tension = 0 and gravity alone provides
centripetal force.)\\
(d) Does the ball have sufficient energy to complete
the full circle?
\hfill\ansref{Ch5-B2}

\vspace{0.4cm}

\textbf{B3.} A petrol-powered water pump on a farm
in Ogun State lifts water from a well $8.0\ \text{m}$
deep and delivers it at $0.025\ \text{m}^3\text{s}^{-1}$
(density of water $= 1000\ \text{kg\,m}^{-3}$,
$g = 10\ \text{m\,s}^{-2}$).\\[4pt]
(a) Calculate the mass of water lifted per second.\\
(b) Calculate the useful power output of the pump.\\
(c) The pump engine produces $3.5\ \text{kW}$ of
input power. Calculate the efficiency.\\
(d) If petrol costs ₦700 per litre and the engine
consumes $0.8\ \text{L\,h}^{-1}$, calculate the
cost of running the pump for 4 hours.\\
(e) How much water (in cubic metres) is lifted in
the 4-hour period?
\hfill\ansref{Ch5-B3}

\vspace{0.4cm}

\textbf{B4.} A $0.20\ \text{kg}$ ball is dropped
from rest onto a spring of spring constant
$800\ \text{N\,m}^{-1}$. The spring is initially
uncompressed, and the ball falls $0.50\ \text{m}$
before first touching the spring.
($g = 10\ \text{m\,s}^{-2}$)\\[4pt]
(a) Find the speed of the ball just before it
touches the spring.\\
(b) Set up an energy equation for the moment of
maximum compression $x$ of the spring (at that
point the ball momentarily stops).\\
(c) Solve for $x$.\\
(d) What is the maximum elastic PE stored in the
spring?
\hfill\ansref{Ch5-B4}

\vspace{0.4cm}

\textbf{B5.} A cyclist and her bicycle have a
combined mass of $75\ \text{kg}$. She rides along
a horizontal road at a constant speed of
$8.0\ \text{m\,s}^{-1}$ against a total resistance
of $40\ \text{N}$.\\[4pt]
(a) Calculate the power she exerts.\\
(b) She then climbs a hill inclined at $5.0°$ to
the horizontal, still at $8.0\ \text{m\,s}^{-1}$
with the same resistance. Find her new power output.\\
(c) At the top of the hill she stops pedalling.
If the same resistance acts, how far does she
travel before stopping? (Use energy methods.)\\
(d) What percentage of her original KE is used
overcoming resistance, and what percentage raises
her height, during the climb of part (b) if she
climbs for $200\ \text{m}$ along the slope?
\hfill\ansref{Ch5-B5}

\end{exercisebox}
